\documentclass[../../main.tex]{subfiles}

\begin{document}

\section{可数集与不可数集}\label{section:ch01-可数集与不可数集}

如同正整数的截是有限集的样板那样，所有正整数的集合 $\mathbb{Z}_{+}$ 就是可数无限集的样板。本节将研究这种集合，还要构造一些既不是有限集也不是可数无限集的集合。这种研究将引导我们去讨论“归纳定义”过程的含义。

\begin{definition}[可数无限集]\label{definition:ch01-可数无限集}
一个集合 $A$ 称为无限的（infinite），如果它不是有限集。一个无限集 $A$ 称为可数无限的（countably infinite），如果存在一个一一对应
\begin{align*}
f: A \longrightarrow \mathbb{Z}_{+}.
\end{align*}
\end{definition}

\begin{example}\label{example:ch01-Z可数}
所有整数的集合 $\mathbb{Z}$ 是可数无限集。容易验证，由
\begin{align*}
f(n) =
\begin{cases}
2n, & \text{如果 } n > 0,\\
-2n + 1, & \text{如果 } n \leqslant 0
\end{cases}
\end{align*}
定义的函数 $f: \mathbb{Z} \longrightarrow \mathbb{Z}_{+}$ 是一个一一对应。
\end{example}

\par
\begin{figure}[htbp]
\centering
\includegraphics[width=0.6\textwidth]{ch01-fig7-1}
\caption{$\mathbb{Z}_{+} \times \mathbb{Z}_{+}$ 的对角枚举}\label{figure:ch01-对角枚举}
\end{figure}

\begin{example}\label{example:ch01-ZxZ可数}
笛卡儿积 $\mathbb{Z}_{+} \times \mathbb{Z}_{+}$ 是一个可数无限集。如果我们用第一象限中的整点表示 $\mathbb{Z}_{+} \times \mathbb{Z}_{+}$ 的元素，那么\cref{figure:ch01-对角枚举}的左边告诉你怎样去“数”这些点。就是说怎样把它们与正整数一一对应起来。当然图不能作为证明，但它给出了证明的启示。首先，作一个一一对应 $f: \mathbb{Z}_{+} \times \mathbb{Z}_{+} \longrightarrow A$，其中 $A$ 是由满足 $y \leqslant x$ 的点 $(x, y)$ 组成的 $\mathbb{Z}_{+} \times \mathbb{Z}_{+}$ 的子集，$f$ 定义为：
\begin{align*}
f(x, y) = (x + y - 1, y).
\end{align*}


然后作 $A$ 与正整数集之间的一个一一对应 $g: A \longrightarrow \mathbb{Z}_{+}$，定义为：
\begin{align*}
g(x, y) = \frac{1}{2}(x - 1)x + y.
\end{align*}
我们把 $f$ 和 $g$ 都是一一对应的证明留给读者完成。后面将给出 $\mathbb{Z}_{+} \times \mathbb{Z}_{+}$ 是可数无限集的另外一个证明。
\end{example}

\begin{definition}[可数集与不可数集]\label{definition:ch01-可数集}
一个集合称为可数的（countable），如果它是有限集或者可数无限集。一个集合不是可数的，就称为不可数的（uncountable）。
\end{definition}

以下定理为我们提供了一个判定集合可数性的常规方法。

\begin{theorem}[可数集的判定]\label{theorem:ch01-可数判定}
设 $B$ 是一个非空集，则下列条件等价：
\begin{enumerate}[(1)]
\item $B$ 是可数集.
\item 存在一个满射 $f: \mathbb{Z}_{+} \longrightarrow B$.
\item 存在一个单射 $g: B \longrightarrow \mathbb{Z}_{+}$.
\end{enumerate}
\end{theorem}

\begin{proof}
(1) $\Longrightarrow$ (2)。设 $B$ 是一个可数集。如果 $B$ 是可数无限集，则根据定义存在一个一一对应 $f: \mathbb{Z}_{+} \longrightarrow B$，从而结论成立。如果 $B$ 是一个有限集，则对于某一个 $n \geqslant 1$，存在一个一一对应 $h: \{1, \cdots, n\} \longrightarrow B$（注意 $B \neq \varnothing$）。可以把 $h$ 扩张为一个满射 $f: \mathbb{Z}_{+} \longrightarrow B$，其定义为：
\begin{align*}
f(i) =
\begin{cases}
h(i), & \text{对于 } 1 \leqslant i \leqslant n,\\
h(1), & \text{对于 } i > n.
\end{cases}
\end{align*}
从而结论成立。

(2) $\Longrightarrow$ (3)。设 $f: \mathbb{Z}_{+} \longrightarrow B$ 是一个满射。$g: B \longrightarrow \mathbb{Z}_{+}$ 定义为：
\begin{align*}
g(b) = f^{-1}(\{b\}) \text{ 的最小元}.
\end{align*}
因为 $f$ 是一个满射，$f^{-1}(\{b\})$ 非空，所以 $g$ 的定义是确切的。如果 $b \neq b'$，集合 $f^{-1}(\{b\})$ 与 $f^{-1}(\{b'\})$ 无交，它们的最小元不同，因此 $g$ 是一个单射。

(3) $\Longrightarrow$ (1)。设 $g: B \longrightarrow \mathbb{Z}_{+}$ 是一个单射，我们要证明 $B$ 是一个可数集。通过改变 $g$ 的值域可以得到 $B$ 与 $\mathbb{Z}_{+}$ 的一个子集之间的一个一一对应。因此只要证明 $\mathbb{Z}_{+}$ 的任意子集是可数的，便可以证明结论。令 $C$ 为 $\mathbb{Z}_{+}$ 的一个子集。

如果 $C$ 是一个有限集，根据定义它是可数的。因此，我们只需证明 $\mathbb{Z}_{+}$ 的无限子集 $C$ 是一个可数无限集。这个论断是容易理解的。事实上，我们可以将 $C$ 的元素排成一个无穷序列，这只需将 $\mathbb{Z}_{+}$ 的元素先按通常的顺序排列，然后再“删除”$\mathbb{Z}_{+}$ 中所有不在 $C$ 中的元素。

以上陈述仅是对证明的轮廓的一个直观说明，严格的证明还须细致的推理分析。我们将作为一个引理予以陈述。
\end{proof}

\begin{lemma}[正整数的无限子集可数]\label{lemma:ch01-正整数无限子集}
设 $C$ 是 $\mathbb{Z}_{+}$ 的一个无限子集，那么 $C$ 是一个可数无限集。
\end{lemma}

\begin{proof}
我们将构造一个一一对应 $h: \mathbb{Z}_{+} \longrightarrow C$。采用归纳法。定义 $h(1)$ 为 $C$ 的最小元，由于 $\mathbb{Z}_{+}$ 的任意非空子集 $C$ 有最小元，所以 $h(1)$ 是存在的。假定 $h(1), \cdots, h(n-1)$ 已有定义，我们定义
\begin{align*}
h(n) = [C - h(\{1, \cdots, n-1\})] \text{ 的最小元}.
\end{align*}
集合 $C - h(\{1, \cdots, n-1\})$ 是非空的，这是因为：如果它是空集的话，$h: \{1, \cdots, n-1\} \longrightarrow C$ 便是一个满射，（根据\cref{corollary:ch01-有限并与积}）$C$ 成为有限集了。所以 $h(n)$ 的定义是确切的。根据归纳法，对于所有 $n \in \mathbb{Z}_{+}$，$h(n)$ 有定义。

容易证明 $h$ 是一个单射。当 $m < n$ 时，注意 $h(m)$ 属于集合 $h(\{1, \cdots, n-1\})$，根据定义 $h(n)$ 不属于这个集合。因此 $h(n) \neq h(m)$。

为了证明 $h$ 是一个满射，令 $c$ 为 $C$ 的任意一个元素，我们证明 $c$ 在 $h$ 的像集中。首先注意到 $h(\mathbb{Z}_{+})$ 不可能被包含在有限集 $\{1, \cdots, c\}$ 中，这是由于 $h(\mathbb{Z}_{+})$ 是无限集（$h$ 是单射）。因此在 $\mathbb{Z}_{+}$ 中有一个元素 $n$，使 $h(n) > c$。令 $m$ 是 $\mathbb{Z}_{+}$ 中满足 $h(m) \geqslant c$ 的最小元，于是对于所有 $i < m$，总有 $h(i) < c$。所以 $c$ 不属于集合 $h(\{1, \cdots, m-1\})$。由于 $h(m)$ 定义为集合 $C - h(\{1, \cdots, m-1\})$ 中的最小元，从而 $h(m) \leqslant c$。两个不等式合起来就得到要证明的 $h(m) = c$。
\end{proof}

在上述定理的证明中，有一处用到了逻辑学的一个原理。当时我们说“应用归纳原理”对所有正整数 $n$ 定义了函数 $h$。你可能以前见过这样的论证，并且当时不曾对它的合法性有所怀疑。在第\cref{section:ch01-整数与实数}的习题中定义 $a^{n}$ 时，我们自己就曾经使用过。

但这里有一个问题。毕竟归纳原理只表明：当 $A$ 是包含正整数集的一个归纳集时，$A = \mathbb{Z}_{+}$。但是，在使用这个原理“根据归纳法”证明一个定理的时候，在一开始就说“设 $A$ 是使得定理成立的所有正整数 $n$ 的集合”，然后回过头来证明 $A$ 为归纳集，从而 $A$ 等于 $\mathbb{Z}_{+}$。

可是，在上面的定理中，我们没有真正用归纳法证明一个定理，而只是用它定义了一些东西。那么我们应该如何着手证明呢？能不能在开始时就说“令 $A$ 是使函数 $h$ 有定义的所有整数 $n$ 的集合”呢？那是不妥当的，因为在证明之初，符号 $h$ 并没有意义。它的意义只是在证明的过程中才出现。因此，先定义一些东西就十分必要了。

这样我们就需要有另一个原理，称之为归纳定义原理（principle of recursive definition）。在前面一个定理的证明中，我们希望能肯定下面的事实：

给定 $\mathbb{Z}_{+}$ 的无限子集 $C$，存在唯一的函数 $h: \mathbb{Z}_{+} \longrightarrow C$，满足条件：
\begin{align}
h(1) = C \text{ 的最小元},\label{align:ch01-归纳公式}
\end{align}
\begin{align*}
h(i) = [C - h(\{1, \cdots, i-1\})] \text{ 的最小元，对于所有 } i > 1.
\end{align*}
公式\eqref{align:ch01-归纳公式}称为 $h$ 的一个归纳公式（recursion formula），它定义了函数 $h$。根据这样的公式给出的定义称为归纳定义（recursive definition）。

在我们试图归纳定义某件事的时候，就会遇到逻辑上的困难。并不是所有归纳公式都有意义。例如，归纳公式
\begin{align*}
h(i) = [C - h(\{1, \cdots, i+1\})] \text{ 的最小元}
\end{align*}
就自相矛盾，虽然 $h(i)$ 必须属于集合 $h(\{1, \cdots, i+1\})$，公式却说它不属于这个集合。下面这个例子也是经典的悖论：

塞维利亚的理发师给塞维利亚的每一个不给自己剃胡子的人剃胡子。那么谁给理发师剃胡子呢？

在这句话里，“理发师”出现了两次，一次在短语“塞维利亚的理发师”中，另一次是作为“塞维利亚人”集合中的一个成员，被理发师剃胡子的人的定义是归纳的。但这个定义是自相矛盾的。

然而，有些归纳公式是有意义的。特别地，有以下原理。

\begin{axiom}[归纳定义原理]\label{axiom:ch01-归纳定义}
设 $A$ 是一个集合。给出一个公式：定义 $h(1)$ 为 $A$ 的唯一一个的元素。对于 $i > 1$ 定义 $h(i)$ 为 $A$ 的这样唯一一个的元素，它仅与 $h$ 在小于 $i$ 的正整数时的值有关，这个公式便决定了唯一的一个函数 $h: \mathbb{Z}_{+} \longrightarrow A$。
\end{axiom}

这就是我们在证明\cref{lemma:ch01-正整数无限子集}时实际用到的原理。如果读者愿意的话，可以直接相信它是对的。当然它也能用归纳原理加以严格证明。下节将给出更精确的描述，并且指出如何予以证明。数学家们很少专门援引这个原理。他们更愿意像在\cref{lemma:ch01-正整数无限子集}的证明中那样，当确实要用到归纳定义原理时，再设法用“归纳原理”来定义一个函数。为了避免过分迂腐，本书也遵从这个惯例。

\begin{corollary}\label{corollary:ch01-可数子集}
可数集的子集是可数的。
\end{corollary}

\begin{proof}
假定 $A \subset B$，其中 $B$ 是一个可数集。那么，存在一个从 $B$ 到 $\mathbb{Z}_{+}$ 的单射 $f$，而 $f$ 在 $A$ 上的限制即为从 $A$ 到 $\mathbb{Z}_{+}$ 的一个单射。
\end{proof}

\begin{corollary}\label{corollary:ch01-ZxZ可数}
集合 $\mathbb{Z}_{+} \times \mathbb{Z}_{+}$ 是可数无限的。
\end{corollary}

\begin{proof}
根据\cref{theorem:ch01-可数判定}，只要构造一个单射 $f: \mathbb{Z}_{+} \times \mathbb{Z}_{+} \longrightarrow \mathbb{Z}_{+}$。将 $f$ 定义为
\begin{align*}
f(n, m) = 2^{n}3^{m}.
\end{align*}
容易验证 $f$ 是一个单射。假定 $2^{n}3^{m} = 2^{p}3^{q}$，当 $n < p$ 时，有 $3^{m} = 2^{p-n}3^{q}$，这与 $3^{m}$ 对于所有 $m$ 都是奇数这一事实相矛盾，因此，$n = p$，从而 $3^{m} = 3^{q}$。如果 $m < q$，则 $1 = 3^{q-m}$，又导致矛盾，因此，$m = q$。
\end{proof}

\begin{example}\label{example:ch01-正有理数可数}
正有理数集 $\mathbb{Q}_{+}$ 是可数无限的。可以用公式
\begin{align*}
g(n, m) = m/n
\end{align*}
定义一个满射 $g: \mathbb{Z}_{+} \times \mathbb{Z}_{+} \longrightarrow \mathbb{Q}_{+}$。由于 $\mathbb{Z}_{+} \times \mathbb{Z}_{+}$ 是一个可数集，所以有一个满射 $f: \mathbb{Z}_{+} \longrightarrow \mathbb{Z}_{+} \times \mathbb{Z}_{+}$。于是复合 $g \circ f: \mathbb{Z}_{+} \longrightarrow \mathbb{Q}_{+}$ 是一个满射。因此，$\mathbb{Q}_{+}$ 是可数集。由于 $\mathbb{Q}_{+}$ 包含着 $\mathbb{Z}_{+}$，它当然也是无限集。有理数集 $\mathbb{Q}$ 是可数无限集的证明留作习题。
\end{example}

\begin{theorem}[可数集的可数并]\label{theorem:ch01-可数并}
可数集的可数并是可数的。
\end{theorem}

\begin{proof}
设 $\{A_{n}\}_{n \in J}$ 为可数集的一个加标族，其指标集 $J$ 为 $\{1, \cdots, N\}$ 或 $\mathbb{Z}_{+}$。为方便起见，还假定每一个集合 $A_{n}$ 非空——这种假定并不影响结论的一般性。

因为每一个 $A_{n}$ 是可数的，故对每一个 $n$ 可以选取一个满射 $f_{n}: \mathbb{Z}_{+} \longrightarrow A_{n}$。类似地，可以选取一个满射 $g: \mathbb{Z}_{+} \longrightarrow J$。用
\begin{align*}
h(k, m) = f_{g(k)}(m)
\end{align*}
定义函数
\begin{align*}
h: \mathbb{Z}_{+} \times \mathbb{Z}_{+} \longrightarrow \bigcup_{n \in J} A_{n}.
\end{align*}
容易验证 $h$ 是一个满射。因为 $\mathbb{Z}_{+} \times \mathbb{Z}_{+}$ 与 $\mathbb{Z}_{+}$ 之间有一个一一对应，根据\cref{theorem:ch01-可数判定}即得出并的可数性。
\end{proof}

\begin{theorem}[可数集的有限积]\label{theorem:ch01-可数积}
可数集的有限积是可数的。
\end{theorem}

\begin{proof}
我们首先证明两个可数集 $A$ 和 $B$ 的积是一个可数集。若 $A$ 或 $B$ 是空集，结论是显然的。若 $A$ 和 $B$ 都不是空集，可选取两个满射 $f: \mathbb{Z}_{+} \longrightarrow A$ 与 $g: \mathbb{Z}_{+} \longrightarrow B$。于是用 $h(n, m) = (f(n), g(m))$ 定义的函数 $h: \mathbb{Z}_{+} \times \mathbb{Z}_{+} \longrightarrow A \times B$ 是一个满射，因此 $A \times B$ 是可数的。

以下对一般情形进行归纳。假定当每一个 $A_{i}$ 可数时，$A_{1} \times \cdots \times A_{n-1}$ 可数，证明对于积 $A_{1} \times \cdots \times A_{n}$ 同样的结论成立。

首先注意，存在一个用
\begin{align*}
g(x_{1}, \cdots, x_{n}) = ((x_{1}, \cdots, x_{n-1}), x_{n})
\end{align*}
定义的一一对应
\begin{align*}
g: A_{1} \times \cdots \times A_{n} \longrightarrow (A_{1} \times \cdots \times A_{n-1}) \times A_{n}.
\end{align*}
因为根据归纳假定 $A_{1} \times \cdots \times A_{n-1}$ 是可数集，根据假定 $A_{n}$ 是可数的，上段已证这样的两个集合的积是可数集，所以 $A_{1} \times \cdots \times A_{n}$ 也是可数的。
\end{proof}

如果能论断可数集的可数积是可数的，那就太好了，可惜这个论断并不成立。

\begin{theorem}[二元素集的 $\omega$ 积不可数]\label{theorem:ch01-二元不可数}
设 $X$ 表示二元素集 $\{0, 1\}$。那么集合 $X^{\omega}$ 不可数。
\end{theorem}

\begin{proof}
我们证明，任意函数
\begin{align*}
g: \mathbb{Z}_{+} \longrightarrow X^{\omega},
\end{align*}
都不是满射。为此，将 $g(n)$ 表示成
\begin{align*}
g(n) = (x_{n1}, x_{n2}, x_{n3}, \cdots, x_{nm}, \cdots),
\end{align*}
其中每一个 $x_{ij}$ 为 $0$ 或 $1$。然后定义 $X^{\omega}$ 的元素 $\boldsymbol{y} = (y_{1}, y_{2}, \cdots, y_{n}, \cdots)$，使得
\begin{align*}
y_{n} =
\begin{cases}
0, & \text{如果 } x_{nn} = 1,\\
1, & \text{如果 } x_{nn} = 0.
\end{cases}
\end{align*}
（如果我们把数 $x_{nn}$ 排成一个矩形方阵，那么其中元素 $x_{nn}$ 总在这个方阵的对角线上，我们选取 $\boldsymbol{y}$ 使其第 $n$ 个坐标不同于对角线上的 $x_{nn}$。）

于是，$\boldsymbol{y}$ 是 $X^{\omega}$ 的一个元素，并且 $\boldsymbol{y}$ 不在 $g$ 的像中。对于给定的 $n$，点 $g(n)$ 与点 $\boldsymbol{y}$ 至少有一个坐标（不妨设它是第 $n$ 个坐标）不相同。这证明 $g$ 不是满射。
\end{proof}

笛卡儿积 $\{0, 1\}^{\omega}$ 是不可数集的一个例子。下面的定理是说明 $\mathcal{P}(\mathbb{Z}_{+})$ 是不可数集的另一个例子。

\begin{theorem}[Cantor 定理]\label{theorem:ch01-Cantor}
设 $A$ 是一个集合。那么不存在单射 $f: \mathcal{P}(A) \longrightarrow A$，也不存在满射 $g: A \longrightarrow \mathcal{P}(A)$。
\end{theorem}

\begin{proof}
一般说来，若 $B$ 非空，则单射 $f: B \longrightarrow C$ 的存在蕴涵着满射 $g: C \longrightarrow B$ 的存在。我们可以对 $f$ 像集中的每一个 $c$ 定义 $g(c) = f^{-1}(c)$，而对于 $C$ 的其余点随意定义它在 $g$ 下的像。

因此，我们只要证明对于任何一个给定的映射 $g: A \longrightarrow \mathcal{P}(A)$，$g$ 不是满射。对于任何一个点 $a \in A$，$a$ 的像 $g(a)$ 为 $A$ 的子集，它可能包含 $a$，也可能不包含 $a$。令 $B$ 为 $A$ 中所有使 $g(a)$ 不包含 $a$ 的点 $a$ 组成的子集，即
\begin{align*}
B = \{a \mid a \in A - g(a)\}.
\end{align*}
$B$ 可能为空集，也可能为整个 $A$，这是没有关系的。可以认定 $B$ 是 $A$ 的一个子集，它不在 $g$ 的像中。假定对于某个 $a_{0} \in A$，$B = g(a_{0})$。试问，$a_{0}$ 是不是一定属于 $B$ 呢？根据 $B$ 的定义，
\begin{align*}
a_{0} \in B \Longleftrightarrow a_{0} \in A - g(a_{0}) \Longleftrightarrow a_{0} \in A - B.
\end{align*}
在任何情况下都导致矛盾。
\end{proof}

现在已经证明了不可数集的存在性。但是还没有提到我们最熟悉的一个不可数集——实数集。你可能以为 $\mathbb{R}$ 的不可数性已经证明过了。如果假定任意实数可以用一个无限十进制数来表示的话（规定在表示式的最后不允许有无限多个 $9$），那么改变一下\cref{theorem:ch01-二元不可数}证明中用过的对角程序，就可以证明实数的不可数性。但是，这个证明在有些方面不十分完善，一个原因是实数的无限十进制表示并不是公理的简单推论，它需要费大力气去证明。另一个原因是 $\mathbb{R}$ 的不可数性实际上并不依赖于 $\mathbb{R}$ 无限十进制展开，也不依赖于 $\mathbb{R}$ 的任何代数性质，它仅与 $\mathbb{R}$ 的序性质有关。我们将晚些时候（第 27 节）用 $\mathbb{R}$ 的序性质来证明 $\mathbb{R}$ 的不可数性。

\subsection*{习题}

\begin{exercise}\label{exercise:ch01-可数集-1}
证明 $\mathbb{Q}$ 是可数无限集。
\end{exercise}

\begin{exercise}\label{exercise:ch01-可数集-2}
证明\cref{example:ch01-Z可数}和\cref{example:ch01-ZxZ可数}中的映射 $f$ 与 $g$ 都是一一对应。
\end{exercise}

\begin{exercise}\label{exercise:ch01-可数集-3}
设 $X$ 为二元集合 $\{0, 1\}$。证明在 $\mathcal{P}(\mathbb{Z}_{+})$ 与笛卡儿积 $X^{\omega}$ 之间有一个一一对应。
\end{exercise}

\begin{exercise}\label{exercise:ch01-可数集-4}
(a) 一个实数 $x$ 称为（在有理数集上）是代数的（algebraic），如果它满足某一个具有有理系数 $a_{i}$ 的正次数的多项式方程
\begin{align*}
x^{n} + a_{n-1}x^{n-1} + \cdots + a_{1}x + a_{0} = 0.
\end{align*}
假定每一个多项式方程仅有有限个根，证明代数数的集合是可数集。

(b) 如果一个实数 $x$ 不是代数的，它称为超越的（transcendental）。在实数集合是不可数的假定下，证明超越数的集合是不可数的。（令人吃惊的是，我们熟悉的超越数仅仅只有两个：e 和 $\pi$。而证明它们是超越数却相当不容易。）
\end{exercise}

\begin{exercise}\label{exercise:ch01-可数集-5}
判定下述每一个集合是不是可数的，并说明理由。
\begin{enumerate}[(a)]
\item 所有函数 $f: \{0, 1\} \longrightarrow \mathbb{Z}_{+}$ 的集合 $A$.
\item 所有函数 $f: \{1, \cdots, n\} \longrightarrow \mathbb{Z}_{+}$ 的集合 $B_{n}$.
\item 集合 $C = \bigcup_{n \in \mathbb{Z}_{+}} B_{n}$.
\item 所有函数 $f: \mathbb{Z}_{+} \longrightarrow \mathbb{Z}_{+}$ 的集合 $D$.
\item 所有函数 $f: \mathbb{Z}_{+} \longrightarrow \{0, 1\}$ 的集合 $E$.
\item 所有“终端为 $0$”的函数 $f: \mathbb{Z}_{+} \longrightarrow \{0, 1\}$ 的集合 $F$. [如果存在正整数 $N$，使得对于所有 $n \geqslant N$，$f(n) = 0$，则称 $f$ 是终端为 $0$（eventually zero）的.]
\item 所有终端为 $1$ 的函数 $f: \mathbb{Z}_{+} \longrightarrow \mathbb{Z}_{+}$ 的集合 $G$.
\item 所有终端为常值的函数 $f: \mathbb{Z}_{+} \longrightarrow \mathbb{Z}_{+}$ 的集合 $H$.
\item $\mathbb{Z}_{+}$ 的所有二元子集的集合 $I$.
\item $\mathbb{Z}_{+}$ 的所有有限子集的集合 $J$.
\end{enumerate}
\end{exercise}

\begin{exercise}\label{exercise:ch01-可数集-6}
如果在集合 $A$ 与 $B$ 之间有一个一一对应，则称 $A$ 和 $B$ 具有相同的基数（same cardinality）。
\begin{enumerate}[(a)]
\item 证明：若 $B \subset A$ 并且存在一个单射
\begin{align*}
f: A \longrightarrow B,
\end{align*}
则 $A$ 和 $B$ 具有相同的基数. [提示：定义 $A_{1} = A$ 和 $B_{1} = B$。对于 $n > 1$，定义 $A_{n} = f(A_{n-1})$ 和 $B_{n} = f(B_{n-1})$。（又是归纳定义！）注意 $A_{1} \supset B_{1} \supset A_{2} \supset B_{2} \supset A_{3} \supset \cdots$。根据
\begin{align*}
h(x) =
\begin{cases}
f(x) & \text{若对于某 } n \text{ 有 } x \in A_{n} - B_{n},\\
x & \text{其他情形}
\end{cases}
\end{align*}
定义函数 $h: A \longrightarrow B$.]
\item 定理（Schroeder-Bernstein 定理） 如果存在单射 $f: A \longrightarrow C$ 及 $g: C \longrightarrow A$，则 $A$ 与 $C$ 具有相同的基数.
\end{enumerate}
\end{exercise}

\begin{exercise}\label{exercise:ch01-可数集-7}
证明习题\cref{exercise:ch01-可数集-5}中的集合 $D$ 与 $E$ 的基数相同。
\end{exercise}

\begin{exercise}\label{exercise:ch01-可数集-8}
设 $X$ 表示二元素集 $\{0, 1\}$，$\mathscr{B}$ 为 $X^{\omega}$ 的所有可数子集的集合。证明 $X^{\omega}$ 与 $\mathscr{B}$ 具有相同的基数。
\end{exercise}

\begin{exercise}\label{exercise:ch01-可数集-9}
(a) 归纳公式
\begin{align}
h(1) = 1,\label{align:ch01-自相矛盾归纳}
\end{align}
\begin{align*}
h(2) = 2,
\end{align*}
\begin{align*}
h(n) = [h(n+1)]^{2} - [h(n-1)]^{2}, \quad \text{对于 } n \geqslant 2
\end{align*}
不符合归纳定义原理。但是可以证明确实存在一个函数 $h: \mathbb{Z}_{+} \longrightarrow \mathbb{R}$，满足这个公式。[提示：改写\eqref{align:ch01-自相矛盾归纳}，使原理适合并且要求 $h$ 为正的。]

(b) 证明 (a) 中的\eqref{align:ch01-自相矛盾归纳}不唯一决定 $h$。[提示：如果 $h$ 是满足\eqref{align:ch01-自相矛盾归纳}的一个正函数，对于 $i \neq 3$，令 $f(i) = h(i)$，而令 $f(3) = -h(3)$。]

(c) 证明：任何函数 $h: \mathbb{Z}_{+} \longrightarrow \mathbb{R}$ 都不满足归纳公式
\begin{align*}
h(1) = 1,
\end{align*}
\begin{align*}
h(2) = 2,
\end{align*}
\begin{align*}
h(n) = [h(n+1)]^{2} + [h(n-1)]^{2}, \quad \text{对于 } n \geqslant 2.
\end{align*}
\end{exercise}

\end{document}